% !TeX program = xelatex
\documentclass[11pt,a4paper]{article}
\usepackage[margin=24mm]{geometry}
\usepackage{amsmath,amssymb,booktabs,array,tikz}
\usetikzlibrary{arrows.meta,positioning}
\usepackage[hidelinks]{hyperref}
\title{The Sixteen Endpoint Classes:\\Definitions, Restricted Proof, and Open Scope}
\author{Internal proof audit}
\date{November 21, 2026}
\begin{document}\maketitle
\noindent\textbf{Status.} A mathematical argument for the specific unit-pair
domain tested by our code, not a theorem for arbitrary connected swept-region
intersections. The main manuscript and implementation are unchanged. Novelty
relative to the literature is not established by this audit.

\section{The actual domain of the classifier}
Let $P_i,P_j$ be nontrivial, simple, oriented paths on an embedded square grid
with $2$ m edge length. Edges join adjacent vertices and are horizontal or
vertical. Each robot has a fixed, axis-aligned $1.5\times1.5$ m square footprint.
Paths may not be left, reversed, or jumped over. Waiting is permitted at path
vertices, and robots remain physically present after reaching their goals.
Motion is continuous between vertices. Rotation leaves the envelope unchanged.

Write $s_i,e_i,s_j,e_j$ for the endpoints. Let
\begin{equation}
K=|P_i|\cap|P_j|,
\end{equation}
where $|P|$ is the geometric centerline image. The tested \emph{unit-pair}
condition is that the common vertex/edge subgraph is connected or empty.
For these simple embedded paths, a nonempty $K$ is a common arc or a singleton.
The code uses the test $|E_K|=|V_K|-1$ when $V_K\ne\varnothing$.

\textbf{Important distinction.} This is not the definition that the two
\emph{swept volumes} have a connected intersection. It is a centerline condition.
The endpoint table's code-level domain must not silently be replaced by an
arbitrary-body geometric domain.

The four bits, in their actual order in \texttt{safe\_order}, are
\begin{equation}
(s_i^K,e_i^K,s_j^K,e_j^K)
=\bigl(\mathbf1_{s_i\in|P_j|},\mathbf1_{e_i\in|P_j|},
\mathbf1_{s_j\in|P_i|},\mathbf1_{e_j\in|P_i|}\bigr).
\label{eq:bits}
\end{equation}
Superscript $K$ distinguishes a Boolean from a position. The code variables
are \texttt{si, ei, sj, ej}; bits refer to membership in the \emph{other full
centerline}, not intersection with an independently chosen local segment.

\section{Why centerline bits also test endpoint sweep interference here}
For a selected grid vertex $q$ not on the other path, the minimum $L_\infty$
distance to that path is at least $2$ m. Two footprints require at most $1.5$ m
of center separation along each coordinate to overlap. Hence
\begin{equation}
q\notin|P_j|\quad\Longleftrightarrow\quad
B_i(q)\cap S_j(P_j)=\varnothing.\label{eq:separation}
\end{equation}
If $q\in|P_j|$, the other robot visits that center, so its full sweep intersects
$B_i(q)$. This proves both directions for the stated grid and envelope.
For denser layouts, wider or rotating bodies, or arbitrary curves,
Equation~\eqref{eq:separation} need not hold. The manuscript's geometric endpoint
flags must then be computed from footprints, not vertex membership alone.

\clearpage
\section{The Boolean formula and all sixteen cases}
The implemented predicate is
\begin{equation}
F_{ij}=(\neg s_j^K\land\neg e_i^K)
\lor(\neg s_i^K\land\neg e_j^K).\label{eq:F}
\end{equation}
The first term certifies $i$ completes before $j$ starts; the second certifies
the reverse. If both hold, the code returns $i$ first, a tie convention only.
Its negation has a useful exact expansion:
\begin{equation}
\neg F_{ij}=
\underbrace{s_i^Ks_j^K}_{A}\ \lor\
\underbrace{e_i^Ke_j^K}_{B}\ \lor\
\underbrace{s_i^Ke_i^K}_{C}\ \lor\
\underbrace{s_j^Ke_j^K}_{D}.\label{eq:blocks}
\end{equation}
Products denote logical AND. The nine negative bit patterns reduce to three
types: start interference (A), terminal interference (B), and blocking by a
contained whole route (C or D). There are four Boolean clauses because the
last type has two symmetric cases, not because there are four distinct types.

\begin{center}
\begin{tabular}{cclll}\toprule
Bits & $F_{ij}$ & Constructive order & Proof category & Saved test count\\\midrule
0000 & 1 & $i$ then $j$ (either) & S & 654,464\\
0001 & 1 & $i$ then $j$ & S & 176,144\\
0010 & 1 & $j$ then $i$ & S & 176,144\\
0011 & 0 & --- & D & 12,656\\
0100 & 1 & $j$ then $i$ & S & 176,144\\
0101 & 0 & --- & B & 113,488\\
0110 & 1 & $j$ then $i$ & S & 113,488\\
0111 & 0 & --- & B, D & 10,128\\
1000 & 1 & $i$ then $j$ & S & 176,144\\
1001 & 1 & $i$ then $j$ & S & 113,488\\
1010 & 0 & --- & A & 42,280\\
1011 & 0 & --- & A, D & 5,064\\
1100 & 0 & --- & C & 12,656\\
1101 & 0 & --- & B, C & 10,128\\
1110 & 0 & --- & A, C & 5,064\\
1111 & 0 & --- & A, B, C, D & 1,616\\\bottomrule
\end{tabular}
\end{center}
S is the serial-execution construction below. A and B are the two-entrance and
two-terminal obstructions. C and D mean one whole route is the common arc.
The counts are from the saved \texttt{max\_edges=5} result, not newly measured.
They total 1,799,096 ordered pairs; 320 representative pairs were cross-checked
with the production solver in that run. Coincident initial positions and
disconnected common paths were excluded by the experiment filters.

\paragraph{Two different statements.}
Algebra alone proves that this table matches Equation~\eqref{eq:F}. It does
\emph{not} prove that the formula matches physical feasibility. The next page
supplies the restricted geometric argument for that additional claim.

\clearpage
\section{Restricted proposition and its proof}
\textbf{Proposition.} Under the assumptions on page 1, the fixed pair admits
a collision-free forward/wait execution with persistent endpoints if and only
if $F_{ij}=1$. When feasible, a whole-route serial execution suffices.
This proposition does not require the grid to be $4\times4$ or the paths to
have at most five edges; those are experimental limits, not proof assumptions.

\subsection*{S: Sufficiency}
If $\neg s_j^K\land\neg e_i^K$, robot $j$ can stay at $s_j$ throughout
robot $i$'s route by Equation~\eqref{eq:separation}. Robot $i$ can then remain
at $e_i$ throughout robot $j$'s route by the same equation. This gives a valid
execution. The other term in Equation~\eqref{eq:F} reverses the two robots.
Any permitted finite segment durations can be used. No precise synchronization
is necessary; vertex point-sweep duplication represents the waiting.

\subsection*{Preliminary facts for necessity}
For a nonempty connected intersection of two simple paths, $K$ is a common
subpath. Each robot visits all of $K$ exactly once. If an endpoint of either
full path belongs to $K$, it must be an endpoint of $K$: it cannot be an interior
point of the common arc, since an interior point has path continuation on both
sides. If both endpoints of $P_i$ are in $K$, then $|P_i|=K$; similarly for $j$.

\textbf{No-passing fact.} Two distinct robot centers constrained to the same
arc cannot exchange their order on that arc without coinciding at some time.
Parameterize the arc by $x\in[0,1]$. If the continuous difference $x_i-x_j$
changes sign while both are on it, it vanishes by the intermediate value
theorem. Coincident centers imply footprint collision. The simple-path
assumption prevents leaving the arc and later re-entering it through a detour.

\subsection*{A: Both starts are in $K$}
If the starts coincide, the initial configuration already collides. Otherwise
they are the two ends of $K$, and each robot is directed into $K$ toward the
other end. Neither can leave $K$ before traversing it. Reaching the opposite
ends would exchange their order. Waiting cannot remove this obstruction, by
the no-passing fact.

\subsection*{B: Both goals are in $K$}
Coincident goals violate persistent final occupancy. If distinct, they are the
two ends of $K$, and each robot enters at the end opposite its goal. If both
enter before either finishes, they must pass on the arc. If one finishes
before the other enters, its permanent footprint blocks that other's entrance.
Neither alternative permits completion.

\subsection*{C and D: One whole route is $K$}
Suppose $|P_i|=K$. If robot $j$ also starts or ends in $K$, case A or B
already applies. Otherwise $j$ must enter one end and leave the other, while
$i$ always remains on $K$. At entry $j$ is behind $i$ in its traversal
direction; to leave the other end it must get ahead. Robot $i$ cannot leave
the arc to let $j$ pass, even after finishing. The no-passing fact rules out
completion. The argument is symmetric for $|P_j|=K$.

Equation~\eqref{eq:blocks} exhausts the negative patterns, proving necessity.
The empty-intersection case has all bits zero and is covered by S. A singleton
intersection in any applicable obstruction gives coincident relevant endpoints;
a nontrivial whole path cannot equal that singleton. $\square$

\clearpage
\small
\section{What this establishes, and what remains open}
\begin{center}\small
\begin{tabular}{p{73mm}p{68mm}}\toprule
Statement & Audit status\\\midrule
The four bits and Boolean table match the code & Direct code/algebra check.\\
Positive rows give a serial completion & Proved by endpoint separation; also
valid with geometric sweep flags for more general bodies.\\
Negative rows rule out completion in the stated grid unit-pair domain &
Restricted proof above, using a single common arc and persistent occupancy.\\
The same necessity result for arbitrary connected sweep intersections &
Not proved; not what the classification code tests.\\
Endpoint table decides multiple common subpaths & False as a general inference;
partial-progress witnesses can violate whole-route serialization.\\
The restricted proposition is new to the literature & Not established.\\
Pair feasibility alone decides an arbitrary multi-robot set & Not established;
this note is a two-robot result.\\\bottomrule
\end{tabular}
\end{center}

\subsection*{How it relates to existing work}
Fixed-path coordination through progress-space obstacles is established.
Leroy, Laumond and Sim\'eon [1] construct a resolution-complete geometric
coordination method using obstacle bounding boxes; its discussion also permits
backward progress, so its admissible-motion model is not identical to our
forward-only model. Ghrist, O'Kane and LaValle
[2] study coordination spaces and Pareto-optimal coordinations on roadmaps.
These are substantially broader settings than a Boolean table for one shared
arc. They are not evidence that this exact endpoint proposition appears
verbatim, but they prevent treating the table as a newly discovered general
coordination framework.

The Segment Graph paper [3, Definition 4] requires paths to avoid other robots'
starts and goals for its deadlock-to-cycle theorem. Its Figure 4 explicitly
shows endpoint blocking. Our table distinguishes several endpoint cases, but
being outside that particular assumption is not proof of novelty: other
coordination methods may already handle them.

\subsection*{Recommended use in our paper}
Keep the general geometric endpoint predicate as a \emph{sufficient} serial
certificate. If the restricted equivalence is useful, state the common-arc and
geometric-separation assumptions next to it and give this short proof. Do not
present finite testing as the proof, or the restricted equivalence as a theorem
about arbitrary mobile-robot sweeps. Assess prior-art equivalence before
advertising the result as a contribution.

\subsection*{Verification record}
This audit reads the classifier's bit order, connected-common-subgraph filter,
independent reachability rules and saved result. The saved counts are reproduced
without launching the 1.8-million-pair experiment again. The Boolean expansion
and all table entries can be checked over 16 assignments; saved representative
routes are rechecked separately. No simulation, viewer or controller change is
part of this note. The proof is an internal argument, not an independent
peer-reviewed validation.

\begin{thebibliography}{9}\small
\bibitem{leroy} S. Leroy, J.-P. Laumond, and T. Sim\'eon, ``Multiple Path
Coordination for Mobile Robots: A Geometric Algorithm,'' IJCAI, 1118--1122,
1999. \url{https://ijcai.org/proceedings/1999/159}.
\bibitem{ghrist} R. Ghrist, J. M. O'Kane, and S. M. LaValle, ``Computing
Pareto Optimal Coordinations on Roadmaps,'' author-hosted manuscript.
\url{https://msl.cs.illinois.edu/~lavalle/papers/GhrOkaLav05.pdf}.
\bibitem{segment} A. Coskun, J. M. O'Kane, and M. Valtorta, ``Deadlock-Free
Online Plan Repair in Multi-robot Coordination with Disturbances,'' FLAIRS,
2021. \url{https://jokane.net/pubs/CosOKaVal21.pdf}.
\end{thebibliography}
\end{document}
